\quellseite{389}
\chapter[§ 27. Holy Scripture]{§ 27. HOLY SCRIPTURE}

*G. v. Zezschwitz ii 2: \emph{Der bibl. Unterricht in der Volksschule}, pp. 188–220. K. Buchrucker, pp. 154–158. E. Sachsse, pp. 334–336. *F. Zange, pp. 73–142. A. Eckert, pp. 33–35; 55–61. O. Baumgarten, pp. 74–80; 87–88. J. Gottschick, pp. 126–141. *J. Berndt, pp. 36–55. *R. Kabisch, pp. 129–156. E. Chr. Achelis, pp. 266–283. A. Rude, pp. 90–99; 126–129. J. Steinbeck, pp. 193–200. E. Kautzsch, \emph{Bibelwissenschaft und Religionsunterricht}, 1900. T. Raymont, \emph{The Use of the Bible in the Education of the Young}, 1911. H. E. Weber, \emph{Historisch-kritische Schriftforschung und Bibelglaube}, ²1914. *Geo. Hodges, \emph{The Training of Children in Religion}, 1917, pp. 167–217. *M. Reu, \emph{Quellen zur Geschichte des kirchlichen Unterrichts, ii: Der Bibelunterricht}, 1906. On Bible Passages: G. W. W. Thilo, \emph{Der Bibelspruch im Dienst des Religionsunterrichts}, 1846. M. Reu, \emph{Quellen etc.}, 1906. Eckhart, \emph{Kurze Erklaerung der wichtigsten Bibelsprueche}, 1887. Peschel, \emph{Ausfuehrliche Erklaerung der wichtigsten Bibelsprueche fuer den Katechismusunterricht}, 1887. L. Schultze, \emph{Katechetische Bausteine}, ⁵1891. *K. Buchrucker, \emph{Der Schriftbeweis im Katechismus}, 1893. Hauffe, \emph{120 Sprueche schulmaessig erklaert}, 1896. *H. Toegel, \emph{Der konkrete Hintergrund zu den 150 Kernspruechen d. relig. Memorierstoffes}, ²1904. H. Toegel, \emph{Die Behandlung der Sprueche} in \emph{Paedag. Studien}, 1904. Th. Meinhold, \emph{Die biblische Grundlage f. d. Katechismusunterricht}. 1907. Zeissig and Fritzsche, \emph{Prakt. Volksschulpaedagogik} (containing seven catecheses on Bible Passages by Toegel), 1907. W. W. Stellhorn, \emph{Schriftbeweis im luth. Katechismus}, 1912. *L. Wessel, \emph{The Proof Texts of the Catechism with a Practical Commentary} in \emph{Theological Quarterly}, 1917 f., book edition 1923. On Bible Reading, Introduction to the Bible, School Bibles and Bible Readers: Boettcher, \emph{Handreichung f. d. Bibellesestunde}, 1894. O. Zuck, \emph{Das Bibellesen im Anschluss an Bibl. Geschichte und Katechismus}, 1895. Melinat, \emph{Das Bibellesen im Volksschulunterricht}, 1898. Habermas, \emph{Handbuch d. Bibellesens u. der Bibelkunde}, 1898. H. and F. Falcke, \emph{Bibellesen}, 1899. J. H. A. Fricke, \emph{Bibelkunde}, i ²1904; ii 1898. J. A. W. Haas, \emph{Bible Literature}, 1903. General Council Graded System, \emph{Bible Outlines, Scripture Quarterly, Lesson Commentary}. H. C. Alleman, \emph{The Book and the Message}, 1914. *M. Reu, \emph{Wartburg Lesson Helps}, course iii: \emph{The Book of Books}, 1926; \emph{The Book of Life}, 2 vols., 1917 ff. John Schaller, \emph{The Book}
\quellseite{390}
\emph{of Books}, 1918. F. Sheatsley, \emph{A Guide to the Study of the Bible}, 1918. R. Hofmann, \emph{Die Schulbibel}, 5¹896. Martin, \emph{Die Schulbibelfrage}, 1888. Baehlnisch, \emph{Ist eine Schulbibel notwendig und wie muss sie beschaffen sein?} 1892. J. Schlier, \emph{Bibel, Schulbibel, Biblisches Lesebuch} in \emph{Neue kirchl. Zeitschr.}, 1904. Behr, \emph{Die Schulbibelfrage}, 1895. Evers, \emph{Die Schulbibelfrage}, 1895. Mueller, \emph{Bibel oder Schulbibel}, 1896. Dix, \emph{Neuere Geschichte der Schulbibel}, 1898. R. Locke, \emph{Vollbibel, Schulbibel, Bibl. Lesebuch}, 1898. Bremische Bibelgesellschaft, \emph{Biblisches Lesebuch oder Schulbibel}, 1894. Schaefer and Krebs, \emph{Bibl. Lesebuch f. d. Schulgebrauch}, 1896. Voelker and Strack, \emph{Biblisches Lesebuch}, 1890. 1898. \emph{Biblisches Lesebuch f. evang. Schulen}, Stuttgart, 1901. Frischi, Schremmer and Holzinger, \emph{Bibl. Lesebuch}, 1909. *Ostermai, Toegel and Neuberg, \emph{Bibl. Lesebuch mit Bildschmuck}, 1911. *Thisselton Mark, \emph{The Young People's Bible}, 2 vols., 1910. *Ella Broadus Robertson, \emph{The Heart of the Bible}, 1911. Kent and Foster, \emph{The Shorter Bible}, 1923. H. Spanuth, \emph{Die Propheten d. A. Bundes. Lebensbilder u. Entwurfe z. unterrichtl. Behandlg.}, 1903. P. Staude, \emph{Der Prophet Jeremiah}; \emph{Praeparationen in darstellender Form}, 1906. Witzmann, \emph{Die unterrichtl. Behandlung d. Psalmen}, 1906. Werkmeister, \emph{20 ausgew. Psalmen; Praeparationen}, 1908. E. Greenwald, \emph{Questions on the Gospels for the Church Year}, 1873. Id., \emph{Questions on the Epistles for the Church Year}, 1874. Sperber, \emph{Erklaerung der Sonn- und Festtagsevangelien}, 31888. R. Kabisch, \emph{Die Episteln u. Evangelien f. Volksschullehrer}, 1897. Lohof, \emph{Kurze Handreichung z. Behandlg. d. Perikopen d. Kirchenjahrs}, 1902.

Biblical History and the Catechism are based exclusively on the Scriptures and in part literally reproduce the Bible text. This might seem to render superfluous further, specific instruction in Holy Scripture. The omission of the Bible from the religious curriculum, however, would not correspond to the supreme importance which Evangelicals attach to that book. As the record and the permanent re-presentation of the divine revelation; as the fountain and norm of all saving knowledge; as the voice of true authority and the book of life for the mature congregation, from which comfort, counsel, incentive, and power are to be drawn every day, the Bible requires a special method of treatment: and instruction in Biblical History and Catechism would be deficient
\quellseite{391}
at a vital point, did they fail to lead the pupil back to Scripture at every step, and thereby prepare him for an independent use of it. If the members of the mature congregation of the present day do not use the Bible with greater diligence, the cause is, in great part, that in their youth they have not been properly introduced to this precious book.

One feature of the instruction of the youth in Holy Scripture is the literal \emph{memorizing of some of the principal parts of Scripture}. Of late, opposition against this practise has arisen, on the plea that it degrades the Bible to the level of a book of oracles; but the advocates of such a view have failed to comprehend one of the peculiar features of Scripture, viz., that fundamental truths have so crystallized in some of its individual passages that the whole Christ, with His holy love and divine grace, shines forth from them. Nor has due account been taken of the need of the ordinary Christian, who lacks the faculty of surveying large complexes, and who is unable so to retain them that they are at his disposal as soon as needed. It is in harmony with his manner of learning a subject and retaining it for future direction and solace to entrust to his heart and memory a treasure of passages in which this or that truth is concentrated as light is in the sun. The history of the Christian Church bears abundant testimony to the fact that just such individual passages have been the stay and comfort of many.

Bible passages are most effectively memorized when associated with instruction in Biblical History or the Catechism; for the Lutheran catechist feels the need of summarizing the fundamental thoughts of the individual Bible narratives in words of Holy Scripture and, likewise, of proving the doctrinal statements of the Catechism by Scripture. The last-named province suggests the necessity for a fund of \emph{dicta probantia}, or proof passages. In the selection of these, the
\quellseite{392}
main requisite is that they are clear and readily intelligible, brief and pithy, trenchant and apt, and also easily memorized and retained. Above all, the connection in which they are found in Scripture must tally with that in which they are employed in catechetical instruction: only thus can they possess the force of proof passages. In connection with some of the chief passages, the catechist will do well to enter more fully upon the biblical context. If a number of passages is used in support of a certain doctrine, they should be so arranged as to correspond to the progress of revelation, e.g., when the Bible proof for the divine sonship of Christ is furnished. Particular attention should be devoted to such passages as embody the consummation and fruitage of whole periods of development. But the criterion for the selection of Bible passages is not merely their evidential value, but also their vital content in the form of comfort and admonition. Words that awaken life and quicken life, words of majesty that engrave themselves, as with style of diamond, upon the heart and conscience of the child, to become a real treasure for all after-life, a stand-by in breath and in death—such words require particular attention. Especially the Psalms and the Prophets, with their abundance of pithy maxims, should freely be drawn upon. It is of the utmost importance that the teacher should guard against the mistake of having too many Bible passages memorized. Where there is no regular parish school, two hundred and fifty passages should be the limit; these should be reviewed with such frequency that they become the children’s never-failing treasure. Only thus time is found to explain them—a duty that must never be set aside, and, also, to have the more gifted children in the course of time memorize some larger sections of Scripture, e.g., Ps. 23; 51; 90; 91; Isa. 9:1–7; 40:26–31; 53:1–12; portions of the Sermon on the Mount and of the last sayings of Christ (such as Matt. 25:31–46);
\quellseite{393}
Phil. 2: 5–11; Rom. 3: 21–28; 5: 1–11; 8: 31–39; 1 Cor. 13; Rev. 3: 17–20.

Another element of instruction in the Holy Scriptures has been, since the sixteenth century, the reading and explanation of the pericopes (usually on Saturdays, to facilitate the understanding of the sermon on the day following). Since instruction in Biblical History has become general, the need of this feature is not so urgent as in the past. Where such instruction is given, the children should look up the pericopes in their Bible, not in their Church Book; the portions hard to understand should be carefully explained; in all cases the main thoughts should be made to stand out; and the whole pericope should be explained to the children in an edifying manner, and impressed upon their minds, in keeping with the respective level of knowledge to which they have attained. But still other features must be observed in the endeavor to introduce the young to Scripture: \emph{Regular reading of the Bible} must take place. This will be most effective where special periods are appointed for it, so that larger sections can be read and studied connectedly. Where that proves impracticable, Bible reading should at least regularly accompany instruction in Biblical History and in the Catechism. It is recommended that, when the creation narrative or the First Article is under consideration, such portions be read as Ps. 104; 121; 91; Job 38:1–40:5; Isa. 40:12 ff.; Ps. 8; 1 Kings 17:1–16; Acts 12:1–24; Gen. 32:10–13; Matt. 6:25–34. In connection with the Conclusion of the Commandments the class might read the narrative of the flood, of the destruction of Sodom and Gomorrah, of the destruction of Jerusalem, of the exile of the children of Israel, but also the narrative of Joseph. In connection with the Second Article the following portions would be appropriate: Luke 2:1–14; John 1:1–18; Isa. 53; Ps. 110; with the Third Article Joel 2:28–32; Acts 2:1–47; John 16:1–15. The Third Chief
\quellseite{394}
Part would receive light from John 17:1–26; the Fourth from 2 Kings 5:1–19; John 3; Acts 8:26–40; Rom. 6:3–11. The Fifth Chief Part suggests the reading of John 6; Rev. 3:14–22; 1 Cor. 11:23–32; and the part dealing with Confession Ps. 51; 32. Where there is a good parish school, time will be afforded in connection with instruction in the Catechism to introduce the pupils to Holy Scripture in the manner here shown; where that opportunity does not exist, larger, connected sections of the Bible ought to be read at least during the last weeks of catechetical instruction. Regular reading and study of the Bible should take place in the Bible class of Sunday Schools and in the devotional meetings of the Young People’s Society.

With the reading of the Bible the more important elements of \emph{Introduction to the Bible} as a whole are readily connected. In Biblical History the sequence of the historical books and their main content are learned; the knowledge of the poetic and prophetic books especially must be imparted by means of Bible reading. The content of the Gospels is imparted in a similar way while the order of the Pauline Epistles can be discussed best in connection with the missionary journeys. Thus the most important phases of introduction to the Bible can easily be taught in connection with the rest of religious instruction. The order of biblical books should not be memorized disconnectedly, but in the connection above suggested. The use of rhymed verses for this purpose can be recommended only upon the condition that better ones than those hitherto used shall have been formed. The children must also necessarily be drilled in finding any one passage at a moment’s notice, they must know something about the technical make-up of the Bible, and they must learn the conventional abbreviations of the biblical books. The proper place for exercises of this sort will be catechetical instruction. As the young begin to mature, they should be
\quellseite{395}
informed concerning the more important sections of the chief books. For the initial stage of Bible introduction the appendix to our \emph{Explanation of Luther’s Catechism} might be acceptable (pp. 147–151). For the confirmed the third course of the \emph{Wartburg Lesson Helps} (\emph{Book of Books}; for the more advanced: \emph{Book of Life}) is intended. The whole process of introduction should glow with respect for the Book and joy in it; and the teacher should at all times aim at equipping the young for personal use of the Bible. Otherwise we need not entertain any hopes of training a generation that delights and takes root in the Bible. It is also recommended to link all truths gleaned from the Scriptures directly with the corresponding portions of the Catechism, because it will help to bring clearness into the child’s religious concept world, cf. § 35.

But is it at all proper to put the Bible into the hands of the children? In view of the frank and unvarnished presentation of sexual matters and sins, this has been looked upon as morally doubtful; for this reason so-called \emph{School-Bibles and Biblical Readers} have been produced. For various reasons, a Biblical Reader is quite desirable (Th. Mark and E. B. Robertson have given us such readers); but even when using the complete Bible, we can, in a measure at least, obviate the difficulty if (1) there is no continuous, but only select, reading, and (2) if we tell our children as soon as the Bible has been put in their hands, that this book pictures man as he is: that even the most heinous sins are there called by name; but that this is not done for our entertainment, but for our humiliation and warning, which may be seen from the fact that, generally, the penalty is pictured immediately after the sin; and that, when we come upon such passages we should read them with holy grief and the sincere prayer that the Father in heaven preserve us from similar transgressions. A determined fight against the carnal appetite of the young should be the chief aim.
